\scanpage{140}
\unsure{transposé|donné par} $K'\xrightarrow{\varphi}\operatorname{End}_{\mathrm Q}(E_0)\xrightarrow{u\leadsto{}^tu}\operatorname{End}_{\mathrm Q}(\check E_0)$. Donnons
nous une polarisation sur $E_0$ (le $\mathrm Q$-motif \deleted{\unsure{ss-jacent}} à $E$).
Alors on trouve un isom $\check E_0\simeq E_0(\unsure{p|\rho})$
($E_0$ de poids $p$), donc \deleted{$K'\longrightarrow\operatorname{End}_{\mathrm Q}(\check E_0)$}
$\operatorname{End}_{\mathrm Q}(\check E_0)\simeq\operatorname{End}(E_0(p))\simeq\operatorname{End}(E_0)$, et on a donc défini
une \uline{involution} \inserted{$u\leadsto u'$} sur $\operatorname{End}_{\unsure{\mathrm Q}}(E_0)$. Le motif dual $(\deleted{\unsure{E_\varphi}}\!E_\varphi)^\vee$
\inserted{de $E_\varphi$ (défini par $\varphi$)} est défini par l'hom $\lambda\leadsto\varphi(\lambda)'\underset{\unsure{0|\circ}}{=}\varphi'(\lambda)$ comme
\[
\boxed{(E_\varphi)^\vee=(E_{\varphi'})(p)}
\]
On n'a pas en général un isom $(E_\varphi)^\vee\simeq E_\varphi(p)$ !

Prenons l'exemple des $K'$-motifs de niveau zéro,
poids $0$. La catégorie est équivalente à la catégorie
des vectoriels de dim finie sur $K'$, \deleted{\illegible{1}}
\deleted{\unsure{en}} à opération \inserted{continue} de $\pi=\operatorname{Gal}(\bar k/k)$. Le dual
correspond à l'opération contragrédiente de $\pi$. Les
$K'$-motifs de rang $1$ \inserted{à isom près} correspondent aux
représentations \deleted{de} \inserted{continues} $\pi\xrightarrow{\chi}K'{}^*$, et on a
$(L_\chi)^\vee=L_{\check\chi}$, où $\check\chi(g)=\chi(g^{-1})$. Donc
on a $L_\chi{}^\vee\simeq L_\chi$ ssi $\chi=\check\chi$ i.e. $\chi$ est
à valeurs dans $\pm1$ !
% archive sheet number: 139
